% Generated draft: requires mathematical and source review.
\documentclass[11pt]{article}
\usepackage[margin=25mm]{geometry}
\usepackage{fontspec}
\setmainfont{Latin Modern Roman}
\usepackage{amsmath,amssymb,mathtools,unicode-math}
\setmathfont{Latin Modern Math}
\usepackage{enumitem,needspace,xurl,hyperref}
\setlength{\parindent}{0pt}
\setlength{\parskip}{5pt}
\setlength{\emergencystretch}{3em}
\newcommand{\N}{\mathbb{N}}
\newcommand{\Z}{\mathbb{Z}}
\newcommand{\Q}{\mathbb{Q}}
\newcommand{\R}{\mathbb{R}}
\newcommand{\C}{\mathbb{C}}
\newcommand{\F}{\mathbb{F}}
\newcommand{\A}{\mathbb{A}}
\newcommand{\PP}{\mathbb{P}}
\newcommand{\E}{\mathbb{E}}
\newcommand{\eps}{\varepsilon}
\newcommand{\dd}{\,\mathrm{d}}
\begin{document}
\section*{UnsolvedMath ID: 20000397}
\textbf{Source problem code:} catalogue entry
\subsection{Definitions and mathematical statement}
Kiran Kedlaya asks to compute monodromy for genus-$g$ curves over $K=\F_q(X)$, where the base curve $X$ need not be rational, and to describe the possible Galois images as $C$ varies in $\mathcal M_g(K)$. The extracted coefficient notation is $\Q_p$. The report instead treats auxiliary-prime cohomology: write $p=\operatorname{char}\F_q$, choose $\ell\ne p$, and put
\[
 V_\ell=H^1_{\mathrm{et}}(C_{\overline K},\Q_\ell),\qquad
 \rho_\ell:G_K\to\operatorname{GL}(V_\ell)
\]
for smooth proper geometrically connected $C/K$. This substitution is explicit; characteristic-$p$ cohomology would require a different precise theory. The report studies Zariski closures, not the entire compact $\ell$-adic image requested literally.

With $\F_q$ the full constant field, define $G_{\mathrm{ar}}$ and $G_{\mathrm{geo}}$ as the Zariski closures of the images of $G_K$ and $G_{K\overline{\F}_q}$ respectively, and superscript $0$ as identity component. Polarization places the geometric group in $\operatorname{Sp}_{2g}$ and the arithmetic group in $\operatorname{GSp}_{2g}$. The report claims that $G_{\mathrm{ar}}/G_{\mathrm{geo}}$ is commutative and $(G_{\mathrm{ar}}^0)^{\mathrm{der}}\subset G_{\mathrm{geo}}^0$.

For an isotrivial curve, meaning one becoming constant after a finite extension, the report gives $G_{\mathrm{geo}}^0=1$. To describe the connected arithmetic group, take Frobenius eigenvalues $\alpha_1,\ldots,\alpha_{2g}$ of a constant model over a finite constant extension. Over an algebraic closure of the coefficient field set
\[
 L_1=\{u\in\Z^{2g}:\prod_i\alpha_i^{u_i}=1\},\qquad
 L_{\mathrm{tor}}=\{u:\prod_i\alpha_i^{u_i}\text{ is a root of unity}\}
 =\operatorname{sat}(L_1).
\]
The full cyclic Frobenius closure has character equations $z^u=1$ for $u\in L_1$; its identity component uses $L_{\mathrm{tor}}$. Replacing Frobenius by a finite power does not change the latter torus. This computes the reported connected arithmetic group, not all possible finite geometric or component data of an isotrivial curve.

In genus one the report gives a two-dimensional arithmetic torus for a constant ordinary elliptic curve and a scalar torus for a constant supersingular curve. In the ordinary case $(G_{\mathrm{ar}}^0\cap\operatorname{Sp}_2)^0$ is a nontrivial one-dimensional torus while $G_{\mathrm{geo}}^0=1$, refuting the proposed general identity between them. These tori need not be split over $\Q_\ell$.

For non-isotrivial elliptic curves, under the report's checked characteristic-greater-than-three Igusa hypotheses, it claims arithmetic $\operatorname{GL}_2$ and geometric $\operatorname{SL}_2$ Zariski monodromy. The exact hypotheses and references are absent. The report does not classify every higher-genus Galois image or answer the original ambiguous $\Q_p$ version.
\subsection{Short English statement}
Compute individual and possible monodromy groups of curves over finite-field function fields. The report separates arithmetic from geometric Zariski monodromy, gives Frobenius-lattice calculations for isotrivial curves and genus-one regimes, and shows that taking the determinant-one part of arithmetic monodromy does not generally recover geometric monodromy. Auxiliary-prime and full-image conventions remain essential.
\subsection{Sources}
\begin{itemize}
\item Local registry: \url{data/unsolvedmath/problems.json}, permanent ID $20000397$, source code \texttt{catalogue entry}; original question and supplied research summary.
\item AIM workshop ``Arithmetic statistics over finite fields and function fields'', Question 19, attributed to Kiran Kedlaya: \url{https://aimath.org/pastworkshops/arithstatffieldproblems.pdf}. Canonical location \texttt{aim-algebraic-number-theory-notes.json}, \texttt{notes[48]}; original remarks and literature field empty. No PDF was opened here.
\item The summary names Igusa for a characteristic-greater-than-three elliptic regime but supplies no full citation or theorem hypotheses; no explicit constant ordinary elliptic equation is included.
\item Supplied machine-generated research attempt 1: partial result, problem open at the run and currently partially solved; rigor metadata ``proved'', novelty confidence low. The non-peer-reviewed monodromy synthesis requires independent verification before citation.
\end{itemize}
\end{document}
